% Chapter Template

% Main chapter title
%\chapter[toc version]{doc version}
\chapter{Design and Implementation}

% Short version of the title for the header
%\chaptermark{version for header}

% Chapter Label
% For referencing this chapter elsewhere, use \ref{ChapterTemplate}
\label{DesignAndImplementation}

This chapter presents the detailed design and implementation of the IRIS system, building upon the system requirements and architectural principles established in Chapter~\ref{Introduction}. To address the challenges identified in the literature review (Chapter~\ref{State of the Art}), a secure livestreaming framework with fog offloading capabilities was developed leveraging the Rust programming language, chosen for its memory safety and concurrency guarantees, which reduces common vulnerabilities such as buffer overflows and data race conditions \cite{jung2017rust, matsakis2014rust}. In comparison to C++, Rust’s ownership model and strict compile-time checks help prevent undefined behavior and improve reliability in complex, high-performance systems \cite{sharma2020comparison}, also all code used in this implementation adheres strictly to safe Rust, avoiding the use of \texttt{unsafe} blocks entirely to maintain guarantees of safety.

\section{System Architecture}

The architecture of the system is grounded in the principles of modularity, service orientation, and distributed computing, drawing from established paradigms such as microservices and fog computing \cite{yu2015survey, bonomi2012fog, vaquero2014finding}. This approach enables the separation of concerns, scalability, and maintainability, which are essential for modern streaming platforms. The system is composed of several distinct components, each fulfilling a specialized role within the overall workflow.

At the ingestion layer, the Restreamer application is employed as a dedicated encoder. Its primary function is to capture video streams from heterogeneous sources and perform protocol-agnostic encoding. This decoupling of encoding from distribution is a recognized best practice in distributed multimedia systems, as it allows for greater flexibility in adapting to varying input formats and network conditions \cite{li2016survey, zhang2017video}. The encoded streams are then forwarded to the core streaming server.

The streaming server, implemented in Rust, serves as the central orchestrator for content distribution. It is responsible for managing client connections, routing streams, and handling protocol translation. The server leverages gRPC for inter-process communication, which provides efficient, low-latency, and language-agnostic messaging \cite{wang2016survey}. Specialized services within the server, such as time synchronization and client reporting, are designed to ensure consistent playback and facilitate adaptive control mechanisms. The strict use of safe Rust throughout the codebase, as recommended in systems programming literature \cite{jung2017rust, matsakis2014rust, sharma2020comparison}, mitigates risks associated with memory safety and concurrency, thereby enhancing reliability and security.

Monitoring and observability are integral to the system’s operation. The streaming server exposes a comprehensive set of metrics via a Prometheus exporter, including retransmission rates, latency, and resource utilization. These metrics are collected and visualized in real time using Grafana, which enables operators to assess system health, diagnose issues, and optimize performance \cite{turnbull2018monitoring, barrett2016monitoring}. This observability framework is consistent with cloud-native and distributed system best practices, where continuous monitoring is critical for maintaining service quality.

Client devices participate in the monitoring process by periodically submitting signed reports to the server. These reports contain performance metrics, counters to prevent replay attacks, and detailed statistics such as retransmission counts and bitmap representations of client status. This mechanism supports robust, data-driven monitoring and enhances the security posture of the system \cite{zhu2013multimedia, li2016survey}.

The architecture is designed to support independent scaling of its components, allowing the encoder, streaming server, and monitoring subsystems to be scaled according to demand. By adhering to open standards and distributed system principles, the system is positioned for future expansion and integration with additional services, maintaining high levels of performance and reliability \cite{yu2015survey, bonomi2012fog, mach2017mobile}.

\begin{figure}[h]
    \centering
    \includegraphics[width=0.95\linewidth]{DesignAndImplementation/global_architecture.pdf}
    \caption{Global architecture of the system, illustrating the main components and their interactions.}
    \label{fig:global-architecture}
\end{figure}

\subsection{Streaming Server}
\section{Protocol Design}

\subsection{Threat Model and Attack Mitigation}
\label{sec:threat-model-mitigation}

In addition to the architectural design, the implementation is organized to facilitate modular development, reproducibility, and ease of deployment. The project directory is structured into distinct folders, each corresponding to a major subsystem or supporting resource. For example, the 	exttt{encoder/} and 	exttt{iris/} directories contain the source code for the encoding and streaming server components, respectively, while configuration files, data, and utility scripts are separated into their own dedicated folders. This organization supports clear separation of concerns and simplifies navigation for developers and reviewers.

To further enhance reproducibility and streamline deployment, containerization is employed throughout the implementation. Dockerfiles are provided for building isolated environments for each major component, and 	exttt{docker-compose.yml} files orchestrate multi-container setups for integration testing and production deployment. This approach ensures that dependencies are managed consistently, reduces the risk of environment-specific issues, and enables rapid setup of the system on different platforms \cite{boettiger2015introduction, merkel2014docker}. By leveraging Docker, the implementation adheres to best practices in modern software engineering, supporting both local development and scalable cloud deployment scenarios.
